% Partie : intuitionistic-logic
% Chapitre : semantics
% Section : introduction

\documentclass[../../../include/open-logic-section]{subfiles}

\begin{document}

\olfileid[fr]{int}{sem}{int}

\olsection{Introduction}

Aucune logique ne peut être décrite de façon satisfaisante sans
sémantique, et la logique intuitionniste ne fait pas exception. Pour
la logique classique, la sémantique fondée sur les !!{valuation}s
s'impose comme référence ; pour la logique intuitionniste, plusieurs
sémantiques se font concurrence. Aucune ne rend compte de la
signification des connecteurs d'une manière qui soit entièrement
satisfaisante du point de vue intuitionniste.

La sémantique fondée sur les !!{relational model}s, semblable à celle
des logiques modales, est peut-être la plus répandue. Elle détermine
dans quels mondes les !!{propositional variable}s sont vraies ; ces
mondes sont reliés par une relation d'accessibilité. Cette relation
est toujours un ordre partiel : elle est réflexive, antisymétrique et
transitive.

Intuitivement, on peut considérer les mondes comme des états de
connaissance ou des « situations où l'on dispose de certaines
justifications ». Un état~$w'$ est accessible depuis~$w$ si, pour ce
que nous savons, $w'$ est un état de connaissance futur possible,
compatible avec ce que l'on sait en~$w$. Une proposition connue ne
cesse pas de l'être : si $!A$ est connue en~$w$ et si~$Rww'$, alors
$!A$ reste connue en~$w'$. La « connaissance » est donc monotone
relativement à la relation d'accessibilité.

Si, par analogie avec la logique épistémique, on dit que $!A$ est
« connue » lorsqu'elle est vraie dans tous les états accessibles,
la réflexivité et la persistance du forçage font ici coïncider cette
condition avec le forçage local. Ainsi, $!A \land !B$ est connue
en~$w$ si, dans tous les états accessibles, $!A$ et~$!B$ le sont.
Puisque $R$ est réflexive et que ce qui est forcé reste forcé dans les
états futurs, $!A \land !B$ est connue en~$w$ si et seulement si $!A$
et~$!B$ sont connues en~$w$. Pour la même raison, $!A \lor !B$ est
connue en~$w$ si et seulement si au moins l'une des deux l'est. Ainsi,
pour $\land$ et $\lor$, les conditions de vérité des connecteurs
coïncident avec celles de la logique classique au monde considéré.

En revanche, la condition de vérité de l'implication diffère de celle
de la logique classique. $!A \lif !B$ est connue en~$w$ si et seulement
si, dans aucun monde~$w'$ tel que~$Rww'$, $!A$ n'est connue sans que
$!B$ le soit aussi. Il ne suffit pas que $!A$ ne soit pas connue ou
que $!B$ soit connue en~$w$. Si nous ne connaissons ni $!A$ ni~$!B$
en~$w$, il peut exister un état futur~$w'$ tel que~$Rww'$ où, en~$w'$,
$!A$ devient connue sans que $!B$ le devienne.

Nous ne connaissons $\lnot !A$ que s'il n'existe aucun état futur
possible dans lequel nous connaissions~$!A$. En effet, si $!A$
pouvait devenir connue, il existerait un tel état où $!A$ serait
connue. Comme nous ne pouvons connaître $\lfalse$, nous y
connaîtrions $!A$ sans connaître~$\lfalse$.

Dans cette interprétation, le principe du tiers exclu échoue :
certaines formules~$!A$ ne sont pas encore connues, mais pourraient
le devenir. Pour une telle !!a{formula}~$!A$, ni $!A$ ni~$\lnot !A$
ne sont connues, donc $!A \lor \lnot !A$ ne l'est pas non plus.
Nous savons cependant, par exemple, que $\lnot(!A
\land \lnot !A)$ : aucun état futur dans lequel $!A$ et~$\lnot !A$
seraient toutes deux connues n'est possible. Cela ne dépend pas du
fait que nous connaissions ou non~$!A$ ou~$\lnot !A$.

Les !!^{relational model}s ne constituent pas l'unique sémantique
possible de la logique intuitionniste. Il existe aussi une sémantique
topologique : les propositions y sont interprétées par des ouverts
d'un espace topologique, et les connecteurs par des opérations sur
ces ensembles (par exemple, $\land$ correspond à l'intersection).

\end{document}
